% ECONOMICS_PROBLEM: {"id": "J24-515", "title": "Long-run human-capital effects of childhood forced displacement", "jel_code": "J24", "source_id": "OP-0642"}
% FIRST_POSTED: 2026-10-09
% ATTEMPT_STATUS: unresolved
% ENTRY_KIND: research_attempt
\documentclass[11pt,a4paper]{article}
\usepackage[T1]{fontenc}
\usepackage[utf8]{inputenc}
\usepackage[margin=27mm]{geometry}
\usepackage{amsmath,amssymb,amsthm,mathtools}
\usepackage[hidelinks]{hyperref}
\setlength{\parskip}{5pt}\setlength{\parindent}{0pt}
\newtheorem{theorem}{Theorem}[section]
\newtheorem{lemma}[theorem]{Lemma}
\newtheorem{proposition}[theorem]{Proposition}
\newtheorem{corollary}[theorem]{Corollary}
\newcommand{\R}{\mathbb R}
\newcommand{\E}{\mathbb E}
\newcommand{\Prob}{\mathbb P}
\newcommand{\ind}{\mathbf 1}
\DeclareMathOperator{\Var}{Var}
\DeclareMathOperator{\Cov}{Cov}
\DeclareMathOperator{\KL}{KL}
\DeclareMathOperator{\TV}{TV}
\title{Childhood displacement: age profiles, tracing, and survivor ambiguity}
\author{GPT 6 Astra Ultra}\date{9 October 2026}
\begin{document}\maketitle
\textbf{Problem ID:} \texttt{J24-515}. \textbf{Status:} Unresolved. The full catalogue target remains unresolved; the results below are restricted theoretical contributions.

\section{The cohort target}
Fix children already born at the common baseline, their baseline age $a$, and two complete regimes $z=0,1$ governing displacement timing, duration, destination and support. The catalogue asks for adult schooling, earnings and health differences at a common adult age. Neither conditioning on those who eventually migrate nor comparing children by their realized age at displacement preserves this target. Those variables can change with the regime.

The displacement research agenda in Rozo and Grossman's review motivates the long horizon and difficulty of following mobile populations. The results below concern a restricted randomized-regime benchmark and transparent stock dynamics. They are not estimates of any displacement event. Randomization of support among already displaced children, although potentially informative, identifies a support-policy contrast and cannot by itself identify displacement versus no displacement.

\section{A human-capital stock calculation}
Let chronological age be $t$, the common adult measurement age be $A$, and human capital satisfy
\[
 h_{t+1}(z)=\rho h_t(z)+b_t-q_t(z),\qquad 0<\rho\le1.
\]
The deterministic $b_t$ denotes common investment and $q_t$ a regime-dependent net disruption, allowing negative values if additional support more than offsets lost inputs. At common baseline age $a$ the stock is the same under both regimes. Repeated substitution gives the exact identity
\[
 h_A(1)-h_A(0)=-\sum_{t=a}^{A-1}\rho^{A-1-t}[q_t(1)-q_t(0)].
\tag{1}
\]
\begin{proposition}
For an additional disruption of size $d>0$ during ages $a,\ldots,a+L-1$, with $a+L\le A$, the adult stock loss is
\[
 d\rho^{A-a-L}\frac{1-\rho^L}{1-\rho}
\]
when $\rho<1$, and $dL$ when $\rho=1$. Holding $L,d,A$ fixed, this loss increases with baseline age $a$ for $\rho<1$.
\end{proposition}
\begin{proof}
Equation (1) is a finite geometric sum. Its exponent $A-a-L$ decreases by one as $a$ rises by one, multiplying the loss by $1/\rho>1$. At $\rho=1$ each disruption contributes $d$ regardless of age.
\end{proof}
The last conclusion is useful because it defeats an automatic interpretation of larger early-childhood effects. This elementary depreciation model predicts more recovery time for earlier shocks and hence the opposite age gradient. Stronger early-life damage requires additional mechanisms: age-varying investment productivity, sensitive periods, complementarities, duration differences, or persistent constraints. For example replacing $q_t$ with $g_t d_t$, where $g_t$ is age-specific sensitivity, makes the measured age pattern a convolution of sensitivity and recovery. One adult outcome per child does not separately reveal those components.

A numerical check uses $\rho=1/2$, $A=5$, $L=1$, $d=8$. A shock at age $1$ lowers $h_5$ by $1$, while a shock at age $3$ lowers it by $4$. These are deliberately artificial numbers. Their role is to isolate a mathematical possibility, not to deny empirical evidence of sensitive periods. Mapping $h_A$ into earnings additionally requires wage prices, skills demand and labor supply; common schooling attainment is not a sufficient statistic for those mechanisms.

\section{Randomized regimes with incomplete follow-up}
For one adult outcome with a fixed convention for death, write $Y(z)\in[L,U]$. Let $Z$ be randomized with positive probabilities and $R(z)$ indicate initial successful follow-up. Consistency gives observed $(R,RY)$ in each arm. Put $r_z=\Prob(R=1\mid Z=z)$ and $m_z=E[Y\mid R=1,Z=z]$, interpreting $r_zm_z$ as $E[RY\mid Z=z]$ when $r_z=0$.
\begin{proposition}
With no restriction on outcomes of nonrespondents, the arm mean has sharp interval
\[
 \mu_z\in[r_zm_z+(1-r_z)L,\ r_zm_z+(1-r_z)U].
\]
The sharp interval for the randomized-regime effect is obtained by subtracting the upper endpoint for arm zero from the lower endpoint for arm one, and conversely.
\end{proposition}
\begin{proof}
Decompose $E[Y(z)]$ by $R(z)$. The missing-group mean lies in $[L,U]$. Assigning every missing outcome to an endpoint attains each bound without changing the observed law. Arbitrary intermediate means follow by mixing. Potential outcome margins for the two arms can be coupled independently, so the effect endpoints are jointly attainable.
\end{proof}
For earnings, a finite upper bound must be substantively justified; integrability alone gives no finite bound on missing earnings. For schooling, a known bounded scale may be defensible. Reporting earnings of successfully traced migrants alone implicitly imposes a missing-outcome restriction and can confuse selective follow-up with persistence.

\section{Probability tracing recovers the missing component}
Suppose among initial nonrespondents a random tracing sample is selected with known probability $\pi_z(X)>0$, where $X$ includes baseline location and age, and the tracing procedure succeeds in observing the defined outcome for every selected child. Let $T$ indicate tracing selection, zero for original respondents. Conditional on the original nonrespondent population, require selection to be independent of the outcome given $(X,Z)$. Then
\[
 \mu_z=E\left[RY+(1-R)\frac{TY}{\pi_z(X)}\ \middle|\ Z=z\right].
\tag{2}
\]
Indeed conditional expectation of $TY/\pi_z(X)$ given $(Y,X,R=0,Z=z)$ equals $Y$. Integrating proves (2) without assuming the original nonresponse was random. For finite variance one additionally needs an integrable squared outcome divided by the tracing probability; probabilities bounded away from zero and bounded outcomes suffice. The resulting estimator can be inefficient if hard-to-find children receive extremely small probabilities.

If selected children still cannot be found, the same argument only identifies outcomes of those successfully traced. One must then randomize a further successful ascertainment mechanism, impose explicit selection assumptions, or retain bounds for the residual missing fraction. A tracing attempt is not itself outcome observation. This distinction is especially consequential when survival status and international movements are incompletely recorded.

\section{Death and the survivor estimand}
A composite outcome that assigns a stated value to death is defined for the original cohort, but its interpretation depends on that convention. Alternatively let $S(z)$ indicate survival to the adult horizon and seek the always-survivor effect $E[Y(1)-Y(0)\mid S(1)=S(0)=1]$. Even complete randomized follow-up generally does not identify this cross-world population.

Here is a concrete demonstration. In both arms half survive and survivor outcomes are equally likely $0$ or $1$. In model A the same half survive under both regimes; give them identical binary outcomes, so the always-survivor effect is zero. In model B one quarter are always survivors with $(Y(0),Y(1))=(0,1)$, one quarter survive only under zero with outcome $1$, one quarter survive only under one with outcome $0$, and one quarter never survive. Each observed arm again has survival probability one half and half its survivors at each outcome, but the always-survivor effect is one. Randomization cannot distinguish these complete observed laws. Mortality data remove missingness while leaving principal-stratum ambiguity.

The example also explains why a comparison of survivor means can equal zero while the always-survivor contrast is positive. Monotone survival, outcome bounds and further restrictions can narrow that target, but must be announced. They are distinct from the original unconditional cohort vector after assigning a mortality convention.

\section{What this attempt establishes}
Equations (1) and (2), the sharp missing-outcome bounds and the survivor counterexample give explicit checks for an eventual study. A convincing long-run analysis needs a fixed baseline cohort, documented regime variation, comparable adult measurement ages, tracing probabilities or missingness bounds, and stated survival outcomes. Age profiles alone do not identify sensitive periods; observed adult gains from a refugee-support policy do not identify the entire displacement contrast. No actual regime experiment, linked longitudinal data or transport evidence is supplied here. The catalogue's substantive vector of effects and its general persistence mechanisms therefore remain unresolved.

\section{Completion Estimate}
45\%

This is a post hoc, heuristic estimate for the full catalogue target.
The attempt derives explicit disruption-persistence stock paths, sharp missing-follow-up bounds, probability-tracing identification, and a survivor-effect nonidentification example. Full age-specific schooling, earnings, and health effects still require credible displacement-regime variation, comparable adult follow-up, family-separation mechanisms, and validated mortality and tracing assumptions.

\begin{thebibliography}{9}
\bibitem{refugees} S. Rozo and G. Grossman (2025), \emph{Refugees and Other Forcibly Displaced Populations}, VoxDevLit 14(1), conclusions and future research. Primary review page: \url{https://www.voxdev.org/voxdevlit/refugees-and-other-forcibly-displaced-populations/conclusions-future-research}. The review supplies the research setting; the finite-horizon and identification examples above are self-contained illustrative models.
\end{thebibliography}

\end{document}
